%%%PAGE 166 rot=0 legibility=good summary=上方便签：赛车过弯最短时间（相切、摩擦力提供向心力）；下方剪贴B5题：导轨上金属杆匀速平动且磁场随时间变化（感生+动生）及解答
%%%BLOCK id=166-01 ch=04 sec=圆周运动·过弯最短时间 kind=problem alt=none cont=none y=0.00-0.30
\begin{problem}
\textbf{力学也有相切}

车利用技术 用最短时间过 \unclear{路宽} $d=10\unit{m}$ % “d=”写在“路宽”二字上方；“路宽”二字潦草
\red{$r_1=\unclear{70}\unit{m}$\quad $\mu=1.25$}
% CHECK: 红字 r_1 之后的数值字形像“70”（有可能是 7.5 或 10；由后面 r'=12.5m、θ=53° 的 3-4-5 三角形反推 r_1 应为 10 或 7.5），照原样存疑
%FIG p166_a 0.148 0.056 0.362 0.182 soft
\notefig[0.3]{figures/p166_a.png}
\begin{solution}
\red{$r'^2=r_1^2+\left[r'-\left(r_1-\dfrac{d}{2}\right)\right]^2$}

\red{$r'=12.5\unit{m}$}

\red{$\mu mg=\dfrac{mv'^2}{r'}$}

\red{$v'=12.5\unit{m/s}$}

\red{$\theta=53^\circ$}

\red{$t=\dfrac{2\theta}{360^\circ}\cdot\dfrac{2\pi r'}{v'}=1.85\unit{s}$}
\end{solution}
\end{problem}
%%%END
%%%BLOCK id=166-02 ch=12 sec=电磁感应·动生与感生(磁场变化+杆切割) kind=problem alt=none cont=none y=0.30-0.77
\begin{problem}
% 剪贴印刷题干。“以恒定速度”与“求在 t=6s”被笔圈出；图中有手写数字 15、9、12、9、15、37°、θ 及“动生”“感生”字样，杆MN后有手画的第二位置竖线（带箭头）
B5、如图所示，两根完全相同的光滑金属导轨 $OP$、$OQ$ 固定在水平桌面上，导轨间的夹角为 $\theta=74^\circ$，导轨单位长度的电阻为 $r_0=0.10\,\Omega/\mathrm{m}$。导轨所在空间有垂直于桌面向下的匀强磁场，且磁场随时间变化 $B=k/t$，其中比例系数 $k=2\unit{T\cdot s}$。将电阻不计的金属杆 $MN$ 放置在水平桌面上，在外力作用下，$t=0$ 时刻开始以恒定速度 $v=2\unit{m/s}$ 从 $O$ 点开始向右平动。已知导轨和金属杆均足够长且与导轨接触良好。求在 $t=6\unit{s}$ 时：(1)回路中感应电动势的大小；(2) $MN$ 所受安培力的大小；(3)外力做功的功率和电路的热功率。
%FIG p166_b 0.52 0.298 0.735 0.46
\notefig[0.4]{figures/p166_b.png}
\begin{solution}
感生 $+$ 动生

$B\downarrow$\qquad $B'(t)=-\dfrac{2}{t^2}<0$

$e=BLv+{-}\dfrac{\Delta B}{\Delta t}S$

$=\dfrac{1}{3}\times18\times2+6=18\unit{V}$\qquad 代入

$r=3\,\Omega$\quad $I=6\unit{A}$\quad $F_A=\dfrac{1}{3}\cdot6\cdot18=36\unit{N}$

$P_{\text{外}}=F_{\text{安}}v=72\unit{W}$

$P_{\text{热}}=I^2r=108\unit{W}$
\end{solution}
\end{problem}
%%%END
%%%PAGEEND 166
